\section{Planificación con Scrum}
\label{sec:scrum}

\subsection{Scrum adaptado a un equipo de una persona}
Scrum define tres responsabilidades (Product Owner, Scrum Master y Developers), tres artefactos (Product Backlog, Sprint Backlog e Incremento) y cinco eventos \cite{scrumguide}. Este laboratorio lo hago yo solo, así que adapté el marco sin cambiar su lógica: el trabajo se ordena por valor, se entrega en iteraciones con una meta, y cada iteración termina con algo que se puede ejecutar y revisar.

\begin{table}[H]
	\centering\small
	\begin{tabularx}{\textwidth}{|p{3.4cm}|X|}
		\hline
		\cab\hd{Elemento de Scrum} & \hd{Cómo se aplicó}\\
		\hline
		Product Owner & Asumo el rol, apoyado en la evidencia de la sección~\ref{sec:contexto} y en la guía del laboratorio, que cumple la función de «cliente». El docente valida en la entrega. \\\hline
		Developers y Scrum Master & Los asumo también. Es una limitación: no hay revisión cruzada de pares. La compensé con pruebas automáticas, trazabilidad y criterios de aceptación escritos antes de programar. \\\hline
		Product Backlog & Tabla~\ref{tab:backlog}: ocho historias con prioridad MoSCoW y puntos. \\\hline
		Sprint Backlog & Tareas de cada historia en las tablas~\ref{tab:sprint1} y \ref{tab:sprint2}. \\\hline
		Incremento & Una versión de la aplicación que se ejecuta y se prueba al cierre de cada sprint. \\\hline
		Sprint Planning & Elegí las historias por prioridad y dependencia; fijé una meta y una capacidad por sprint. \\\hline
		Daily Scrum & Se reemplaza por el registro de cambios del repositorio (anexo~\ref{anx:git}): cada commit dice qué se avanzó. \\\hline
		Sprint Review & Revisión de capturas y de resultados de pruebas contra los criterios de aceptación (sección~\ref{sec:resultados}). \\\hline
		Sprint Retrospective & Subsección~\ref{sec:retro}. \\\hline
	\end{tabularx}
	\caption{Elementos de Scrum y su adaptación al trabajo individual.}
	\label{tab:scrumadap}
\end{table}

\onefig{img/d5_scrum.png}{1.0}{Flujo de trabajo: de la evidencia al Product Backlog, dos sprints con su revisión y la entrega.}{fig:scrum}

\subsection{Meta de producto y metas de sprint}
\begin{description}[leftmargin=0pt]
	\item[Meta de producto.] Que una persona matriculada en la UNSA pueda asegurar su cupo de laboratorio desde cualquier lugar y dentro de su turno, obtener ayuda sin hacer fila y no perder su lugar si el servicio falla.
	\item[Meta del Sprint 1.] El estudiante sabe cuándo le toca matricularse, entiende por qué una acción está bloqueada y puede pedir atención sin fila (HU-08, HU-12, HU-09; 13 puntos).
	\item[Meta del Sprint 2.] La matrícula no se pierde si el servicio cae, queda con un comprobante verificable y las pantallas nuevas se leen con texto ampliado (HU-11, HU-10, HU-13; 13 puntos).
\end{description}

Fijé una capacidad de 13 puntos por sprint como límite de trabajo en curso. Es una decisión de planificación mía; no es una velocidad medida y no la presento como tal. HU-11 (la más grande, 8 puntos) va en el Sprint 2 porque reutiliza las validaciones de turno del Sprint 1 y el comprobante de HU-10.

\textbf{Nota de honestidad sobre los sprints.} Los dos sprints son iteraciones lógicas, ordenadas por prioridad y dependencia; no los medí en días ni semanas, y el repositorio los registra juntos en un solo commit de código (anexo~\ref{anx:git}). Por eso no presento una gráfica de avance ni una velocidad: no tengo datos para sostenerlas.

\subsection{Sprint 1: turnos, explicación y atención sin fila}
\begin{table}[H]
	\centering\footnotesize
	\begin{tabularx}{\textwidth}{|p{1.1cm}|X|p{4.6cm}|}
		\hline
		\cab\hd{HU} & \hd{Tareas} & \hd{Archivos principales}\\
		\hline
		HU-08 & Modelo de turno con estados; cuatro turnos de demostración; regla que bloquea la matrícula fuera de turno; selector de ciclo en el inicio de sesión; aviso de turno en Cursos; lista de turnos en Ayuda y en Operación. & \texttt{models.dart}, \texttt{datos\_demo.dart}, \texttt{estado.dart}, \texttt{login\_screen.dart}, \texttt{turno\_widgets.dart} \\\hline
		HU-12 & Función que calcula el motivo del bloqueo; estados del botón del grupo; preguntas frecuentes. & \texttt{estado.dart}, \texttt{grupo\_detalle\_screen.dart}, \texttt{ayuda\_screen.dart} \\\hline
		HU-09 & Modelo de número de atención; solicitar, cancelar y atender; tarjeta con posición y espera; cola del administrador; avisos. & \texttt{models.dart}, \texttt{estado.dart}, \texttt{ayuda\_screen.dart}, \texttt{operacion\_screen.dart} \\\hline
	\end{tabularx}
	\caption{Sprint Backlog del Sprint 1 (13 puntos).}
	\label{tab:sprint1}
\end{table}

\subsection{Sprint 2: contingencia, comprobante y legibilidad}
\begin{table}[H]
	\centering\footnotesize
	\begin{tabularx}{\textwidth}{|p{1.1cm}|X|p{4.6cm}|}
		\hline
		\cab\hd{HU} & \hd{Tareas} & \hd{Archivos principales}\\
		\hline
		HU-11 & Interruptor de caída en Operación; banner de servicio; solicitudes pendientes con código; procesamiento por orden de llegada con motivo de rechazo; sección de pendientes en el horario. & \texttt{estado.dart}, \texttt{turno\_widgets.dart}, \texttt{operacion\_screen.dart}, \texttt{horario\_screen.dart}, \texttt{grupo\_detalle\_screen.dart} \\\hline
		HU-10 & Código de comprobante en cada matrícula; fila «Comprobante» y botón para copiarlo. & \texttt{models.dart}, \texttt{estado.dart}, \texttt{confirmacion\_screen.dart} \\\hline
		HU-13 & Etiqueta semántica del número; pruebas con texto al 150\,\%; corrección del contraste del verde. & \texttt{ayuda\_screen.dart}, \texttt{lab05\_widget\_test.dart}, \texttt{comunes.dart} \\\hline
	\end{tabularx}
	\caption{Sprint Backlog del Sprint 2 (13 puntos).}
	\label{tab:sprint2}
\end{table}

\subsection{Definición de terminado}
Una historia está terminada cuando cumple todos los puntos de la tabla~\ref{tab:dod}. Indico cuáles verifiqué por revisión y cuáles quedan pendientes de la ejecución.

\begin{table}[H]
	\centering\footnotesize
	\begin{tabularx}{\textwidth}{|p{0.7cm}|X|p{3.4cm}|}
		\hline
		\cab\hd{N.°} & \hd{Condición} & \hd{Estado}\\
		\hline
		1 & Cumple todos sus criterios de aceptación (tabla~\ref{tab:criterios}) con pruebas automáticas definidas. & Pruebas escritas; \textbf{resultado pendiente} \\\hline
		2 & \texttt{flutter analyze} no reporta errores. & \textbf{Pendiente} de ejecución \\\hline
		3 & Se ve bien a 360\,dp (teléfono) y a 900\,dp (tableta). & Capturas c30 y c16/c17; \textbf{pendiente} de generarlas \\\hline
		4 & No se desborda con el texto al 150\,\%. & Pruebas escritas; \textbf{pendiente} \\\hline
		5 & El texto de las etiquetas cumple contraste 4,5:1. & Cumple (tabla~\ref{tab:contraste}) \\\hline
		6 & No usa paquetes externos (solo el SDK). & Cumple (\texttt{pubspec.yaml}) \\\hline
		7 & Tiene al menos una captura en el recorrido. & Capturas definidas; \textbf{pendiente} \\\hline
		8 & Está registrada en el repositorio con un mensaje claro. & Cumple (anexo~\ref{anx:git}) \\\hline
	\end{tabularx}
	\caption{Definición de terminado (DoD).}
	\label{tab:dod}
\end{table}

\subsection{Eventos y retrospectiva}
\label{sec:retro}
\paragraph{Revisión de los sprints.} La revisión consiste en ejecutar \texttt{flutter test} y comparar cada captura con su criterio de aceptación. Esa ejecución se hace en mi equipo; hasta que lo haga no afirmo que los criterios se cumplen.

\pendiente{Revisión del Sprint 1 y del Sprint 2: indicar qué criterios pasaron y cuáles fallaron, con los resultados de la sección~\ref{sec:resultados}.}

\paragraph{Retrospectiva (preliminar).} Lo que sé hasta ahora, sin esperar a los resultados:

\begin{table}[H]
	\centering\small
	\begin{tabularx}{\textwidth}{|p{3.4cm}|X|}
		\hline
		\cab\hd{Qué funcionó} & \hd{Evidencia}\\
		\hline
		Partir de la evidencia & Tres de las cuatro funciones principales (turnos escalonados, protocolo de contingencia y mejor orientación) coinciden con medidas que el director de Servicios Académicos nombró como pendientes \cite{elpueblo25}; no las inventé. \\\hline
		Reloj inyectable & Los turnos dependen de la hora; con un reloj fijo las pruebas y las capturas son repetibles y los textos como «faltan 2 d 0 h» no cambian. \\\hline
		Revisar el propio diseño & Calcular el contraste de las etiquetas reveló que el verde heredado del Laboratorio~4 no cumplía; se corrigió. \\\hline
		\cab\hd{Qué mejorar} & \hd{Acción}\\
		\hline
		Un commit por historia & El registro de cambios quedó con un solo commit de código para los dos sprints, por lo que no sirve como Daily Scrum. Acción: en los próximos laboratorios, un commit por historia terminada. \\\hline
		Verificar más temprano & Ejecutar \texttt{flutter analyze} y \texttt{flutter test} al final de cada sprint y no solo al cierre del laboratorio (el script \texttt{ejecutar.ps1} lo hace en un paso). \\\hline
		Datos propios de usuarios & La evidencia es documental. Con más tiempo, una encuesta breve a compañeros validaría los hallazgos H1 a H8 y ajustaría el backlog. \\\hline
		Trabajo en solitario & Sin revisión de pares, un error de criterio solo lo detectan las pruebas. Acción: pedir una revisión corta del informe y del código a un compañero antes de entregar. \\\hline
	\end{tabularx}
	\caption{Retrospectiva preliminar del trabajo.}
	\label{tab:retro}
\end{table}
